\SourceLeaf{29}
\fancyhead[L]{\small ALLGEMEINES GRUNDGESETZ DER MECHANIK.}
\pdfbookmark[0]{Über ein neues allgemeines Grundgesetz der Mechanik}{mechanik}
\begin{center}
ÜBER EIN NEUES\par
\bigskip
ALLGEMEINES GRUNDGESETZ DER MECHANIK.\par
\medskip
\rule{.16\textwidth}{.3pt}
\end{center}
\expandafter\def\csname HistoricalNoteText2\endcsname{Es sei mir jedoch hier die Bemerkung erlaubt, dass ich die Art, wie ein anderer grosser Geometer versucht hat, \textsc{Huyghens} Gesetz für die ausserordentliche Brechung des Lichts in Krystallen von doppelter Brechung, vermittelst des Grundsatzes der kleinsten Wirkung zu beweisen, nicht befriedigend finde. In der That ist die Zulässigkeit dieses Grundsatzes wesentlich von dem der Erhaltung der lebendigen Kräfte abhängig, nach welchem die Geschwindigkeiten der bewegten materiellen Punkte bloss durch ihre Plätze bedingt werden, ohne dass die Richtung der Bewegung Einfluss darauf haben kann, was doch in dem erwähnten Versuch vorausgesetzt wird. Es scheint mir, dass im Emanationssystem alle Bemühungen, die Erscheinungen der doppelten Brechung an die allgemeinen dynamischen Gesetze anzuknüpfen, so lange erfolglos bleiben müssen, als man die Lichttheilchen bloss wie Punkte betrachtet.}
\par
Bekanntlich verwandelt das Princip der virtuellen Geschwindigkeiten die ganze Statik in eine mathematische Aufgabe, und durch \textsc{Dalembert}’s Princip für die Dynamik ist diese wiederum auf die Statik zurückgeführt. Es liegt daher in der Natur der Sache, dass es kein neues Grundprincip für die Bewegungs- und Gleichgewichts-Lehre geben \textit{kann}, welches der Materie nach nicht in jenen beiden schon enthalten und aus ihnen abzuleiten wäre. Inzwischen scheint doch wegen dieses Umstandes noch nicht jedes neue Princip werthlos zu werden. Es wird allezeit interessant und lehrreich bleiben, den Naturgesetzen einen neuen vortheilhaften Gesichtspunkt abzugewinnen, sei es, dass man aus demselben diese oder jene einzelne Aufgabe leichter auflösen könne, oder dass sich aus ihm eine besondere Angemessenheit offenbare. Der grosse Geometer, der das Gebäude der Mechanik auf dem Grunde des Princips der virtuellen Geschwindigkeiten auf eine so glänzende Art aufgeführt hat, hat es nicht verschmäht, \textsc{Maupertuis} Princip der kleinsten Wirkung zu grösserer Bestimmtheit und Allgemeinheit zu erheben, ein Princip, dessen man sich zuweilen mit vielem Vortheil bedienen kann\HistoricalNote{2}.
